\subsection{Confirmation and spillovers in the school cohorts}\label{app:sae-fresh}
We retained the procedures in Appendix~\ref{app:sae}, including every grade, the program-specific realized lottery, entitlement floor and continuity option outside the five discretionary slots. The model uses the 2019 fitted coefficients and 2019 program frequencies in every test cycle. The new regular-round archives supply the same fields used in 2020--21. Application years 2022 and 2023 correspond to admission years 2023 and 2024. The recorded initial course can be compared with menu membership even though the allocation under a changed display is a counterfactual. The sign criterion was specified separately for each year before obtaining these releases. No applicant-level independence or between-year inferential interval is assumed. The two cohorts contain 520,420 and 492,888 officially assigned applicants, respectively. Outside outcomes, including students with no discretionary choice, count toward the placement denominator.

\begin{table}[t]\centering\small
\caption{Previously unexamined school application cycles. Rank-informed minus rank-blind at five discretionary choices. A recorded offer counts as covered if present in the display. The last columns partition the placement change by whether the pupil would receive a place under full-list pooled DA.}
\label{tab:sae-later}
\begin{tabular}{@{}lrrrrr@{}}\toprule
Year & Applicants & $\Delta$ coverage & $\Delta$ placed & $\Delta$ retained & $\Delta$ entrants\\\midrule
2022 & 570,891 & $+1,143$ & $-468$ & $-110$ & $-358$\\
2023 & 536,353 & $+1,350$ & $-517$ & $-117$ & $-400$\\\bottomrule
\end{tabular}\end{table}

\paragraph{Whose list changed?}
The B1 field \textit{prioritario} identifies families with an administrative priority designation. It supplies neither measured household income nor a randomized group assignment. We computed all four combinations of which group's long lists use the rank-informed score. When \emph{only other applicants} switch, every designated pupil has exactly the same menu as before. Among designated pupils in 2022, 127 lost a placement and 67 gained one; in 2023, 125 lost and 60 gained. In both years this is an externality on people who never received the new ranking. Under universal use, the designated group's net placement change per hundred applicants was $-0.050$ and $-0.064$ points, versus $-0.126$ and $-0.148$ among other applicants. Standardizing both groups to the pooled grade, discretionary list-length band and entitlement distribution gives $-0.064$ and $-0.077$ versus $-0.103$ and $-0.126$. All strata for 2022/23 contain both groups. Group composition explains part of the raw gap. These descriptive contrasts do not isolate a causal effect of priority status.

\paragraph{Representing reserved places.}
The primary one-pool replay omits the allocation of earmarked vacancies. As a sensitivity, each program first retains returning pupils, then uses up to its published \textit{vacantes prioritarios} places for remaining designated pupils, allowing any unfilled reserve to revert to the open pool. Each step uses the original coarse priority and realized lottery. This simplified reserve-first rule omits other reserve types, dynamic family links and the ministry's full sequence of assignments. The full-list exact program/outside agreement rises slightly from 95.385\% to 95.414\% in 2022 and from 95.785\% to 95.800\% in 2023. At five choices the ranking switch reduces placements by 469 and 516 under the sensitivity, close to the pooled losses of 468 and 517. Official-offer coverage does not depend on the subsequent allocation rule. The archive gives gross flows, rank changes for those remaining matched, and editing-group decompositions for each of the four years.

\paragraph{Offers families actually accepted.}
As a follow-up we restricted the administrative coverage target to applicants whose parent or guardian explicitly accepted the initial offer, including those who accepted while waiting for a more preferred course. The 2022 and 2023 cohorts have 382,638 and 367,749 such records. Rank-informed menus covered 1,000 and 1,244 more of their accepted initial courses than rank-blind menus. Thus 2,244 of the overall 2,493 net additional covered initial offers belong to explicit acceptances. We did not infer acceptance for nonrespondents or observe what families would have accepted under a shortened menu.
